\documentclass{article}
\usepackage{amsmath, amssymb, amsthm}
\usepackage{hyperref}
\usepackage{geometry}
\geometry{margin=1in}

\title{Erd\H{o}s Problem \#653}
\date{Last updated: 2025-08-31}
\author{Source: erdosproblems.com}

\begin{document}
\maketitle

\noindent\textbf{Status:} OPEN \quad \textbf{No prize}

\noindent\textbf{Tags:} geometry, distances

\medskip

\noindent\textbf{Problem Statement:}

Let $x_1,\ldots,x_n\in \mathbb{R}^2$ and let $R(x_i)=\#\{ \lvert x_j-x_i\rvert : j\neq i\}$, where the points are ordered such that\[R(x_1)\leq \cdots \leq R(x_n).\]Let $g(n)$ be the maximum number of distinct values the $R(x_i)$ can take. Is it true that $g(n) \geq (1-o(1))n$?




Erd\H{o}s and Fishburn proved $g(n)&#62;\frac{3}{8}n$ and Csizmadia proved $g(n)&#62;\frac{7}{10}n$. Both groups proved $g(n) &lt; n-cn^{2/3}$ for some constant $c&gt;0$.

\noindent\textbf{Related OEIS sequences:} possible


\bigskip
\noindent\small{Source: \url{https://www.erdosproblems.com/653}}
\end{document}
